\section{Comparison contract and completed matrix}
\label{sec:experimental-methodology}
The primary comparisons are DG, all-cell OFDG, and OFDG--KXRCF. Here DG means
no oscillation filter; it still receives the common Euler positivity treatment.
OEDG is included on selected affine shock cases. Within each problem, all
methods share the mesh, degree, flux, CFL setting, final time, and boundary data.
RK4 is used throughout. Scalar CFL is 0.15 and Euler CFL is 0.30; the timestep
also depends on degree, element size, and the current wave-speed estimate.
OFDG is applied after complete steps and OEDG at RK stages. Consequently, this
is a comparison of these configured algorithms, not an isolated sensor test.
KXRCF uses its default threshold of one.

\begin{table}[htbp]
\centering\small
\begin{tabular}{@{}p{.33\textwidth}p{.47\textwidth}r@{}}
\toprule Problem & Degree and cells & $T$\\\midrule
Smooth periodic advection & P1--P3; 32, 64, 128, 256 & 1.1\\
Piecewise advection & P2; 128, 256 & 1.1\\
Smooth Burgers & P2; 32, 64, 128 & 0.6\\
Post-shock Burgers & P2; 128, 256 & 2.2\\
Lax & P1/192, P2/128, P3/96 & 1.3\\
Shu--Osher & P2; 200, 400 & 1.8\\
Straight/curved 2D advection & degree 2; $8^2,16^2,32^2$ & 1.0\\
2D Euler Riemann & degree 2; $32^2,64^2$ & 0.25\\\bottomrule
\end{tabular}
\caption{Completed preview matrix. P$k$ denotes degree $k$ in 1D. The 2D
quadrilateral spaces are tensor-product $Q^k$, also called P$k$ in driver
metadata. OEDG accompanies the finer piecewise, post-shock Burgers, and
Shu--Osher runs, plus P2 Lax.}
\label{tab:preview-matrix}
\end{table}

The study completed 103 solver runs, including three temporal-sensitivity runs,
and four independent reference solves. Numerical work took 7.35 minutes on an
eight-processor laptop, sequentially using one rank for 1D and two for the
larger 2D cases. The sampled peak process-tree memory was 118 MiB. All rows
completed; no failed or omitted row is hidden. These single-run timings describe
the resource budget, not comparative performance.

\subsection{Measurements and reference limitations}
For smooth solutions, the principal error is the quadrature evaluation of
$E_h=\|u-u_h\|_{L^2(\Omega)}$. The observed refinement rate is
\begin{equation}
 r=\log(E_h/E_{h/2})/\log 2.\label{eq:observed-rate}
\end{equation}
The reported extrema of discontinuous profiles use 21 samples per 1D cell;
they are not certified extrema of every polynomial. Curves break at element
interfaces, preserving DG jumps. Physical cell averages are exported separately
and are explicitly identified where plotted. A common physical $y$ coordinate
is used for all 2D line cuts; sampling neighbouring nodal rows would compare
different locations.

Euler reference curves come from a separate finite-volume WENO5 implementation
with SSPRK3 time integration. Each of Lax and Shu--Osher has 2,048- and
4,096-cell references. For sensitivity, adjacent fine cell averages are averaged
to the coarse grid before differencing. The fine curve is a numerical reference,
not an exact solution; the maximum differences in \cref{tab:reference-sensitivity}
are substantial near steep features. Small visual separations must therefore
not be interpreted as proven accuracy rankings.

Only periodic scalar runs are used for the unforced conservation-drift claim:
the largest recorded drift is $3.99\times10^{-14}$. For nonperiodic Euler problems,
total content can change through the boundary. A conservation error would
require the time-integrated numerical boundary flux, which this preview does
not supply. Filter preservation of a cell mean is a separate algebraic property.

The approved preview manifest and retained commands define the configurations.
This dataset is separate from both the historical quick workflow and the
unexecuted full-study proposal; see \cref{app:study-design}.
